\section{Introduction}

It is important to design provably safe controllers for autonomous systems. 
Hamilton-Jacobi (HJ) reachability analysis provides a powerful framework to design such controllers for general nonlinear dynamical systems \citep{lygeros2004reachability, mitchell2005time}.
In reachability analysis, safety is characterized by the system's \textit{Backward Reachable Tube (BRT)}.
This is the set of states from which trajectories will eventually reach a given target set despite the best control effort.
Thus, if the target set represents undesirable states, the BRT represents unsafe states and should be avoided. 
Along with the BRT, reachability analysis provides a safety controller to keep the system outside the BRT. 

Traditionally, the BRT computation in HJ reachability is formulated as an optimal control problem.
The BRT can then be obtained as a sub-zero level solution of the corresponding value function.
Obtaining the value function requires solving a partial differential equation (PDE) over a state-space grid, resulting in an exponentially scaling computation complexity with the number of states \citep{bansal2017hamilton}.
To overcome this challenge, a variety of solutions have been proposed that trade off between the class of dynamics
they can handle, the approximation quality of the BRT,
and the required computation.
These include specialized methods for linear and affine dynamics \citep{10.1007/3-540-64358-3_38, Frehse2011, Kurzhanski00, Kurzhanski02, Maidens13, girard2005reachability, althoff2010computing, bak2019numerical, Nilsson2016}, polynomial dynamics \citep{doi:10.1177/0278364914528059, majumdar2017funnel, Dreossi16, henrion2014convex}, monotonic dynamics \citep{coogan2015efficient}, and convex dynamics \citep{chow2017algorithm}
(see \citet{bansal2017hamilton, bansal2021deepreach} for a survey).

Owing to the success of deep learning, there has also been a surge of interest in approximating high-dimensional BRTs \citep{rubies2019classification, fisac2019bridging, djeridane2006neural, niarchos2006neural, darbon2020overcoming} and optimal controllers \citep{onken2022neural} through deep neural networks (DNNs).
Building upon this line of work, \citet{bansal2021deepreach} have proposed DeepReach -- a toolbox that leverages recent advances in neural implicit representations and neural PDE solvers to compute a value function and a safety controller for high-dimensional systems.
Compared to the aforementioned methods, DeepReach can handle general nonlinear dynamics, the presence of exogenous disturbances, as well as state and input constraints during the BRT computation.
Consequently, methods for verifying neural reachable tubes have been proposed. For example, \citet{lin2023generating} propose an iterative scenario-based method \citep{campi2009scenario} to recover probabilistically safe reachable tubes from DeepReach solutions up to a desired confidence level and bound on violation rate. Unfortunately, the method does not allow an after-the-fact risk-return trade-off, and as a result, it is highly sensitive to outlier errors in the learned solutions. This can lead to highly conservative reachable tubes and a severe loss of recovery in the case of stringent safety requirements, as we demonstrate in our case studies.

In this work, we propose two different verification methods, one based on robust scenario optimization and the other based on conformal prediction, to provide probabilistic safety guarantees for neural reachable tubes.
Both methods are resilient to the outlier errors in neural reachable tubes and automatically trade-off the strength of the probabilistic safety guarantees based on the outlier rate.
The proposed methods can evaluate any candidate tube and are not restricted to a specific class of system dynamics or value functions.
We further prove that these seemingly different verification methods naturally reduce to one another, providing a unifying viewpoint for uncertainty quantification (typical use case of conformal prediction) and error optimization (typical use case of scenario optimization) in neural reachable tubes.
Based on these insights, we propose an outlier-adjusted verification approach that can recover a greater safe volume from a neural reachable tube by harnessing information about the distribution of error in the learned solution.
To summarize, the key contributions of this paper are:
\begin{itemize}
    \item probabilistic safety verification methods for neural reachable tubes that enable a direct trade-off between resilience and the probabilistic strength of safety,
    \item a proof that split conformal prediction reduces to a scenario-based approach in general, demonstrating a strong relationship between the two highly related but disparate fields,
    \item an outlier-adjusted verification approach that recovers greater safe volumes from tubes, and
    \item a demonstration of the proposed approaches for the high-dimensional problems of multi-vehicle collision avoidance and rocket landing with no-go zones.
\end{itemize}